% paper/sections/theorem4_algorithm.tex -- P2-05, plan/03-phase2-theory.md.
%
% Standalone section fragment, \input{} by paper/main.tex once P6-01 builds
% the skeleton. Same preamble requirements as the other theoremN files.
% Depends on setup.tex, theorem1_bound.tex, theorem2_lower_bound.tex,
% theorem3_identifiability.tex's notation and machinery. This section
% DEFINES the algorithm P3-03 implements -- P3 should treat this
% definition as the spec, not re-derive its own.

\section{Theorem 4: correctness of Extrapolation-Track-and-Stop}
\label{sec:theorem4}

\subsection{The algorithm}
\label{sec:theorem4-algorithm}

\textbf{Inputs:} $\delta \in (0, 1/2)$, an upper bound $\eta_k \ge \sqrt{\sigma^2_{\mathrm{extrap},k}}$
on each arm's bias magnitude at the current design shape (estimated via P1-06's
debiased estimator, Remark~\ref{rem:sigma2-extrap-estimator} -- see
Section~\ref{sec:theorem4-eta} for why this must be an input, not a discovered
quantity), a design shape $\pi$ (the relative weighting across scales that
Theorem~\ref{thm:lower-bound}'s $T^\star(\nu)$ program optimizes; supplied by P3-02's
solver, plugged in with the current $\hat\theta_k$'s).

\textbf{Tracking (round $t$):} pull $(\arg\min_{k,s} N_k(s,t) / \pi(k,s))$ -- the
most-under-sampled (arm, scale) pair relative to the target allocation $\pi$, the
standard "C-tracking" rule (Garivier \& Kaufmann, 2016) applied to $\pi$ instead of a
per-arm allocation. \emph{Cited, standard}: this specific tracking rule and its
property (empirical allocation converges to $\pi$) are unchanged from the classical
literature; nothing here is specific to the parametric setting.

\textbf{Stopping and recommendation (round $t$):} let $\hat{k}(t) = \arg\max_k
\hat\mu_k(s^\star; t)$ (current leading arm), and for each $k \ne \hat{k}(t)$:
\[
  \widehat{\Delta}_k(t) = \hat\mu_{\hat{k}(t)}(s^\star;t) - \hat\mu_k(s^\star;t), \qquad
  c_k(t, \delta) = \sqrt{2\big(v_{\hat{k}(t)}(t) + v_k(t)\big) \log(1/\beta(t,\delta))}
\]
for an anytime-valid threshold $\beta(t,\delta)$ (e.g. the standard
$\beta(t,\delta) = \delta / (t(t+1))$ union-bound-over-time construction, or a
law-of-the-iterated-logarithm-type bound -- \emph{cited, standard}: this piece is
unchanged from classical sequential testing and not re-derived here). Define
\[
  \textsc{Certified}_k(t) = \big\{ \widehat{\Delta}_k(t) - \eta_{\hat{k}(t)} - \eta_k
    \;>\; c_k(t,\delta) \big\}, \qquad
  \textsc{Abstain}_k(t) = \big\{ c_k(t,\delta) \le \epsilon_0 \ \text{ and }\
    \widehat{\Delta}_k(t) - \eta_{\hat{k}(t)} - \eta_k \;\le\; c_k(t,\delta) \big\}
\]
for a small fixed tolerance $\epsilon_0 > 0$ (a practical "close enough" cutoff for
"the confidence radius has stopped shrinking usefully"; any fixed $\epsilon_0$ works for
Theorem~\ref{thm:algorithm-correctness}'s asymptotic claims, since the abstention event
below is driven by $\eta$, not $\epsilon_0$, as $t \to \infty$).
\begin{itemize}
  \item \textbf{Stop and output $\hat{k}(t)$} if $\textsc{Certified}_k(t)$ holds for
    every $k \ne \hat{k}(t)$.
  \item \textbf{Abstain, output "cannot certify; recommend $s_{\mathrm{rec}}$"} if
    $\textsc{Abstain}_k(t)$ holds for any $k \ne \hat{k}(t)$, where $s_{\mathrm{rec}}$ is
    the single-scale (Corollary~\ref{cor:single-scale}) recommendation at the largest
    scale actually observed.
  \item Otherwise, continue tracking.
\end{itemize}

\subsection{Why a bias floor, not a growing threshold}
\label{sec:theorem4-eta}

\begin{remark}[The naive Track-and-Stop never stops here]
Classical Track-and-Stop's GLR statistic is (up to constants) $\widehat{\Delta}_k(t)^2 /
(2(v_{\hat{k}}(t)+v_k(t)))$, compared against a threshold $\sim \log(t/\delta)$ that
\emph{grows} with $t$ -- correct when $v_k(t) \to 0$ and $\widehat\Delta_k(t)$ converges
to the \emph{true} gap $\Delta_k > 0$, since then the statistic diverges and eventually
crosses any growing threshold. Here, $\widehat\Delta_k(t) \to \Delta_k +
\mathrm{bias}(D_k)$ (Corollary~\ref{cor:pairwise}), not $\Delta_k$ itself -- if
$|\mathrm{bias}(D_k)|$ is large enough to matter, the naive statistic converges to a
\emph{finite} limit and a growing threshold is eventually never crossed: the algorithm
runs forever. This is exactly why route (b) (condition on a known/estimated $\eta$,
rather than route (a), inflating the confidence level) is used:
$\textsc{Certified}_k(t)$'s threshold does \emph{not} grow -- it compares a
\emph{shrinking} $c_k(t,\delta)$ against a \emph{fixed} bias floor
$\eta_{\hat{k}}+\eta_k$, so the algorithm correctly recognizes when further compute
cannot help, rather than chasing a threshold it can never reach.
\end{remark}

\begin{theorem}[$\delta$-correctness, conditional on $\eta$]
\label{thm:algorithm-correctness}
If $\eta_k \ge \sqrt{\sigma^2_{\mathrm{extrap},k}}$ for every $k$ (the inputs are valid
upper bounds), then whenever the algorithm stops and outputs $\hat{k}$,
$\mathbb{P}[\hat{k} \ne k^\star] \le \delta$.
\end{theorem}

\begin{proof}
Mechanical, directly from Theorem~\ref{thm:main-bound}'s corrected bound and the
anytime-valid threshold construction: at the (random) stopping time $\tau$,
$\textsc{Certified}_k(\tau)$ holding for every $k \ne \hat{k}(\tau)$ means
$\widehat\Delta_k(\tau) - \eta_{\hat{k}(\tau)} - \eta_k > c_k(\tau,\delta)$, i.e. (by the
same worst-case-bias argument as Theorem~\ref{thm:main-bound}'s Step 3-4, now with $\eta$
in place of $\sqrt{\sigma^2_{\mathrm{extrap}}}$ since $\eta$ is assumed to \emph{upper
bound} it) the event $\{\hat{k}(\tau) \ne k^\star\}$ implies some $c_k(\tau,\delta)$-scale
deviation event whose probability, union-bounded over both $k$ and every possible
stopping time $\tau$ (the anytime-valid $\beta(t,\delta)$ construction's entire purpose),
is at most $\delta$. Not re-derived in full here since it is a direct combination of two
already-proven pieces (Theorem~\ref{thm:main-bound}'s corrected worst-case bound, and the
standard anytime-valid-threshold union-bound-over-time argument, \emph{cited} from the
classical sequential-testing literature) rather than new content.
\end{proof}

\subsection{Asymptotic optimality in the solvable regime}
\label{sec:theorem4-optimality}

\begin{theorem}[Asymptotic optimality]
\label{thm:asymptotic-optimality}
In the solvable regime (Theorem~\ref{thm:impossibility}'s hypothesis fails --
informally, $\eta_{k^\star} + \eta_k < \Delta_k$ for every $k$, i.e. the true gaps
exceed the bias floor), tracking allocation $\pi = w^\star(\nu)$ (Theorem
\ref{thm:lower-bound}'s optimizer) gives
\[
  \frac{\mathbb{E}_\nu[C(\tau)]}{\log(1/\delta)} \;\longrightarrow\; T^\star(\nu)
  \qquad \text{as } \delta \to 0.
\]
\end{theorem}

\begin{proof}[Proof sketch]
\needshuman{The upper-bound half is the standard Track-and-Stop asymptotic-optimality
argument (Garivier \& Kaufmann, 2016, Theorem 14 and its proof technique): C-tracking
drives the empirical allocation to $\pi$, the anytime-valid threshold is tight enough
that stopping occurs as soon as $\textsc{Certified}$ triggers (no meaningful
over-shoot), and in the solvable regime the bias floor $\eta_{\hat{k}}+\eta_k$ is
eventually dominated by $\Delta_k$, so the stopping condition asymptotically reduces to
the classical (bias-free) one -- adapting that proof's tracking-convergence and
stopping-time lemmas to this parametric setting (rather than a per-arm empirical mean)
is mechanical in structure but has not been written out line-by-line here, and is
exactly the kind of "mechanical adaptation" that still needs a careful pass to confirm
no parametric-specific step silently breaks a classical lemma's hypothesis (e.g. the
tracking-convergence lemma's proof typically uses per-arm sample-mean concentration
directly; here it needs Theorem~\ref{thm:main-bound} Step 1's M-estimator concentration
instead, the same nonlinear-$g$ caveat flagged throughout this document). The matching
lower bound is Theorem~\ref{thm:lower-bound} directly, no new argument.}
\end{proof}

\subsection{Graceful abstention in the impossible regime}
\label{sec:theorem4-abstention}

\begin{theorem}[Abstention triggers with high probability in the impossible regime]
\label{thm:abstention}
Suppose the current design shape places nonzero weight on scales
$s < s^\star$ only (Theorem~\ref{thm:impossibility}'s setting) and the true instance
is one of a pair $\nu_1, \nu_2$ satisfying that theorem's hypothesis with
$\eta_k \ge |\mathrm{bias}(D_k)|$ under either instance (i.e. the assumed $\eta$'s
correctly bound \emph{this specific} construction's perturbation). Then, as $t \to
\infty$ along the tracking rule, $\mathbb{P}[\textsc{Abstain}_k(t) \text{ eventually
holds for the confusable } k] \to 1$; the algorithm does not run forever without
resolution.
\end{theorem}

\begin{proof}
Since $\nu_1,\nu_2$ agree on every scale $s<s^\star$ (Theorem~\ref{thm:impossibility}'s
hypothesis), $\hat\theta_{k}(t)$'s distribution -- and hence
$\widehat\Delta_k(t)$'s -- is identical whether the true instance is $\nu_1$ or $\nu_2$
(Theorem~\ref{thm:impossibility}'s proof, same argument). C-tracking under either
instance drives $v_{\hat{k}}(t) + v_k(t) \to 0$ (Theorem~\ref{thm:lower-bound} Step 2,
unaffected by which instance is true, since it only depends on the design), so
$c_k(t,\delta) \to 0$ almost surely, eventually falling below any fixed $\epsilon_0$.
By Theorem~\ref{thm:main-bound} Step 2 (delta-method consistency),
$\widehat\Delta_k(t) \to D_k^\dagger$ a.s. (the population pairwise difference under
the true instance). Theorem~\ref{thm:impossibility}'s construction (Section
\ref{sec:theorem2-construction}) is built exactly so that
$|D_k^\dagger| \le \eta_{\hat{k}}+\eta_k$ under (at least) one of $\nu_1,\nu_2$ -- by
hypothesis here, this holds for the true instance. Hence $\widehat\Delta_k(t) -
\eta_{\hat{k}}-\eta_k \to D_k^\dagger - \eta_{\hat{k}}-\eta_k \le 0 \le c_k(t,\delta)$
eventually (since $c_k(t,\delta) \ge 0$ always), so both conjuncts of
$\textsc{Abstain}_k(t)$ hold from some finite (random) time onward almost surely --
convergence a.s. implies convergence in probability, giving the stated limit. Mechanical,
given the already-proven pieces cited.
\end{proof}

\begin{remark}[This is the practically valuable part]
Theorem~\ref{thm:abstention} is what makes the algorithm a genuine \emph{diagnostic}, not
just a best-arm-identification routine that happens to sometimes be wrong: rather than
silently running past its useful horizon (or silently returning an answer with a
falsely-inflated confidence), it detects -- from the same quantities
(Theorem~\ref{thm:main-bound}'s bias/variance decomposition) already computed for every
other purpose in this project -- exactly the situation Theorem~\ref{thm:impossibility}
proves is hopeless at this design, and falls back to the one baseline
(Corollary~\ref{cor:single-scale}) that needs no extrapolation assumption at all.
\end{remark}

\subsection{Numerical certificate}
\label{sec:theorem4-certificate}

\texttt{tests/theory/test\_theorem4.py} implements the \textsc{Certified}/\textsc{Abstain}
stopping-and-abstention logic exactly as specified above (not the full adaptive
C-tracking algorithm, which needs P3-02's $T^\star$ solver and is P3-03's job to build
and P3-04's to simulate end-to-end) against simulated compute-increasing sequences at a
\emph{fixed} design shape, using the same linear-in-theta simulation machinery as
Theorem~\ref{thm:main-bound}'s certificate. Checks: (a) in a solvable instance (true
gap exceeds the assumed $\eta$ budget), the rule stops and selects $k^\star$ correctly
in $\ge 1-\delta$ of repeated simulated runs; (b) in an instance built via
Theorem~\ref{thm:impossibility}'s exact construction (Section
\ref{sec:theorem2-construction}), the rule abstains (rather than stopping incorrectly,
or never terminating within a generous compute cap) in effectively every run; (c) the
algorithm never both stops \emph{and} abstains, and never falsely abstains when the
true gap comfortably exceeds the $\eta$ budget.
